
\section{\BookRecordIndex{T0143}格罗莫夫--威滕不变量的公理}\label{sec:7-5}
\zhindex{geluomofuweitengbubianliangdegongli}{格罗莫夫--威滕不变量的公理}

让我们回到本章开头的启发式讨论。那里把同调类
$A\in H_2(M;\mathbb Z)$ 和整数 $k\ge3$ 的\BookRecordIndex{T0143}格罗莫夫--威滕不变量描述为同态
\[
  \operatorname{GW}_{A,k}^M:
  H^*(M)^{\otimes k}\otimes H_*(\overline{\mathcal M}_{0,k})
  \longrightarrow\mathbb Q,
\]
其形式定义为
\[
  \operatorname{GW}_{A,k}^M(a_1,\ldots,a_k;\beta)
  :=\int_{\overline{\mathcal M}_{0,k}(A;J)}
  \operatorname{ev}_1^*a_1\smile\cdots\smile
  \operatorname{ev}_k^*a_k\smile\pi^*\operatorname{PD}(\beta).
\]
回忆，$\operatorname{ev}_i:\overline{\mathcal M}_{0,k}(A;J)\to M$ 是在第 $i$ 个\BookRecordIndex{T0066}标记点处的
\BookRecordIndex{T0283}求值映射，而
\[
  \pi:\overline{\mathcal M}_{0,k}(A;J)
  \longrightarrow\overline{\mathcal M}_{0,k}
\]
是\BookRecordIndex{T0377}遗忘映射，它把每个\BookRecordIndex{T0340}稳定映射送到其底\BookRecordIndex{T0072}层约化\BookRecordIndex{T0339}稳定曲线。在
$\overline{\mathcal M}_{0,k}(A;J)$ 上积分，应理解为把上同调类在\BookRecordIndex{T0246}模空间的\BookRecordIndex{T0369}虚基本循环上
取值。

在这里讨论的一般性下，即使把\BookRecordIndex{T0143}格罗莫夫--威滕不变量取在整（上）同调类上，也必须允许
它取有理数值。不过，如果 $(M,\omega)$ \BookRecordIndex{T0058}半正，那么亏格零不变量在整同调类上取整数值。
高亏格时会出现有理数值，因为当 $g>0$ 时曲线\BookRecordIndex{T0246}模空间 $\overline{\mathcal M}_{g,k}$ 只是
一个轨形；非\BookRecordIndex{T0058}半正情形也会出现有理数值，因为此时以\BookRecordIndex{T0370}虚模循环定义计数不变量时，
\BookRecordIndex{T0113}多重覆盖曲线会起作用。虽然本书并不以这种形式建立理论，考察这一表述下\BookRecordIndex{T0143}格罗莫夫--威滕
不变量的形式性质仍然很有意思。附录~D 将证明，亏格零\BookRecordIndex{T0339}稳定曲线的
Grothendieck--Knudsen \BookRecordIndex{T0246}模空间 $\overline{\mathcal M}_{0,k}$ 是光滑流形。对每棵稳定\BookRecordIndex{T0084}带标号树
$T=(T,E,\Lambda)$，以 $T$ 为模型的\BookRecordIndex{T0339}稳定曲线之集的闭包构成光滑子流形
$\overline{\mathcal M}_{0,T}$；由 Keel \srcite{208} 的结果，所有这类子流形共同生成
$\overline{\mathcal M}_{0,k}$ 的同调。（见定理 D.4.7 和 D.6.4；那里提到的子流形 $V_S$
由定理 D.4.7(iv) 都具有 $\overline{\mathcal M}_{0,T}$ 的形式。）

正如本章导言所指出的，$\operatorname{GW}_{A,k}^M$ 在几何上应理解为如下代数计数：计数
$M$ 中表示类 $A$、带 $k$ 个\BookRecordIndex{T0066}标记点的\BookRecordIndex{T0340}稳定映射，它们在各\BookRecordIndex{T0066}标记点处通过与 $a_i$ Poincaré
对偶的循环 $X_i\subset M$，而且其约化\BookRecordIndex{T0339}稳定曲线属于代表类
$\beta$ 的循环 $Y\subset\overline{\mathcal M}_{0,k}$。在第~\ref{sec:7-1}、\ref{sec:7-3}
节中，我们已经在 $(M,\omega)$ \BookRecordIndex{T0058}半正，并且 $\beta$ 或者由某棵 $k$-\BookRecordIndex{T0084}带标号树 $T$ 对应的
子流形 $\overline{\mathcal M}_{0,T}\subset\overline{\mathcal M}_{0,k}$ 表示，或者由\BookRecordIndex{T0150}固定
标记点 $\mathbf w:=\{w_i\}_{i\in I}$ 所得的子流形 $Y_{k,I}$ 表示的假设下，严格定义了
这些不变量。

Kontsevich 和 Manin 在 \srcite{216} 中列出了\BookRecordIndex{T0143}格罗莫夫--威滕不变量的如下公理。他们考察
所有亏格的不变量；这里专门讨论亏格零，并且先假设 $k\ge3$（$0\le k\le2$ 的情形见
注~\ref{rem:7-5-2}）。如定理 D.4.2 所述，我们用\BookRecordIndex{T0185}交比函数
$\mathbf z\mapsto w_{ijk\ell}(\mathbf z)$ 把 $\overline{\mathcal M}_{0,k}$ 与它在
$(S^2)^N$ 中的像 $\overline M_{0,k}$ 等同，因此将其中的元素记作元组
$\{w_{ijk\ell}\}$。

\paragraph{（有效性）}
若 $\omega(A)<0$，则 $\operatorname{GW}_{A,k}^M=0$。

\paragraph{（对称性）}
对每个置换 $\sigma\in S_k$，
\[
  \operatorname{GW}_{A,k}^M
  (a_{\sigma(1)},\ldots,a_{\sigma(k)};\sigma_*\beta)
  =\varepsilon(\sigma;a)\,
  \operatorname{GW}_{A,k}^M(a_1,\ldots,a_k;\beta).
\]
这里
\[
  \varepsilon(\sigma;a)
  :=(-1)^{\#\{i<j\mid \sigma(i)>\sigma(j),\,
                  \deg(a_i)\deg(a_j)\in2\mathbb Z+1\}}
\]
表示在奇数次数类上诱导的置换之符号；$\sigma_*\beta$ 表示
$\beta\in H_*(\overline{\mathcal M}_{0,k})$ 在如下微分同胚下的上推：
\[
  \{w_{i_0i_1i_2i_3}\}
  \longmapsto
  \sigma_*\{w_{i_0i_1i_2i_3}\}
  :=\{w_{\sigma(i_0)\sigma(i_1)\sigma(i_2)\sigma(i_3)}\}.
\]

\paragraph{（分次）}
若 $\operatorname{GW}_{A,k}^M(a_1,\ldots,a_k;\beta)\ne0$，则
\[
  \sum_{i=1}^k\deg(a_i)-\deg(\beta)=2n+2c_1(A).
\]

\paragraph{（同调）}
对每个 $A\in H_2(M;\mathbb Z)$ 和每个整数 $k\ge3$，都存在同调类
\[
  \sigma_{A,k}\in
  H_{2n+2c_1(A)+2k-6}
  \bigl(M^k\times\overline{\mathcal M}_{0,k}\bigr),
\]
使得
\[
  \operatorname{GW}_{A,k}^M(a_1,\ldots,a_k;\beta)
  =\bigl\langle
    \pi_1^*a_1\smile\cdots\smile\pi_k^*a_k
      \smile\pi_0^*\operatorname{PD}(\beta),\,
    \sigma_{A,k}\bigr\rangle.
\]
其中 $\pi_i:M^k\times\overline{\mathcal M}_{0,k}\to M$ 是到第 $i$ 个因子的投影，
$\pi_0:M^k\times\overline{\mathcal M}_{0,k}\to\overline{\mathcal M}_{0,k}$ 是到最后一个因子的
投影。

\paragraph{（基本类）}
令
\[
  \pi_{0,k}:\overline{\mathcal M}_{0,k}
  \longrightarrow\overline{\mathcal M}_{0,k-1}
\]
表示遗忘最后一个\BookRecordIndex{T0066}标记点的映射，并令
$1:=\operatorname{PD}([M])\in H^0(M)$。若 $(A,k)\ne(0,3)$，则
\[
  \operatorname{GW}_{A,k}^M(a_1,\ldots,a_{k-1},1;\beta)
  =\operatorname{GW}_{A,k-1}^M
    (a_1,\ldots,a_{k-1};(\pi_{0,k})_*\beta).
\]
特别地，当 $\beta=[\overline{\mathcal M}_{0,k}]$ 时，这个不变量为零。

\paragraph{（除子）}
若 $(A,k)\ne(0,3)$ 且 $\deg(a_k)=2$，则
\[
  \operatorname{GW}_{A,k}^M
  \bigl(a_1,\ldots,a_k;
    \operatorname{PD}(\pi_{0,k}^*\operatorname{PD}(\beta))\bigr)
  =\operatorname{GW}_{A,k-1}^M(a_1,\ldots,a_{k-1};\beta)
    \int_Aa_k.
\]

\paragraph{（零类）}
若 $A=0$，则当 $\deg(\beta)>0$ 时
$\operatorname{GW}_{0,k}^M(a_1,\ldots,a_k;\beta)=0$，并且
\[
  \operatorname{GW}_{0,k}^M(a_1,\ldots,a_k;[\mathrm{pt}])
  =\int_Ma_1\smile\cdots\smile a_k.
\]

\paragraph{（分裂）}
固定 $H^*(M)$ 的一组基 $e_0,\ldots,e_N$，令
\[
  g_{\nu\mu}:=\int_Me_\nu\smile e_\mu,
\]
并用 $g^{\nu\mu}$ 表示逆矩阵。再固定指标集的分拆
$\{1,\ldots,k\}=S_0\sqcup S_1$，其中
$k_i:=\#S_i\ge2$（$i=0,1$），并令
\[
  \phi_S:\overline{\mathcal M}_{0,k_0+1}
  \times\overline{\mathcal M}_{0,k_1+1}
  \longrightarrow\overline{\mathcal M}_{0,k}
\]
为如下典范映射：它把第一条\BookRecordIndex{T0339}稳定曲线的最后一个\BookRecordIndex{T0066}标记点与第二条\BookRecordIndex{T0339}稳定曲线的第一个\BookRecordIndex{T0066}标记点
\BookRecordIndex{T0254}粘合。其余指标有唯一的次序，使相对次序保持不变，且
$\overline{\mathcal M}_{0,k_0+1}$ 中的前 $k_0$ 个点映到以 $S_0$ 为指标的点，
$\overline{\mathcal M}_{0,k_1+1}$ 中的后 $k_1$ 个点映到以 $S_1$ 为指标的点。固定
\[
  \beta_0\in H_*(\overline{\mathcal M}_{0,k_0+1}),\qquad
  \beta_1\in H_*(\overline{\mathcal M}_{0,k_1+1}).
\]
则
\begin{align*}
 &\operatorname{GW}_{A,k}^M
   (a_1,\ldots,a_k;\phi_{S*}(\beta_0\otimes\beta_1))\\
 &\quad=\varepsilon(S,a)
   \sum_{A_0+A_1=A}
   \operatorname{GW}_{A_0,k_0+1}^M
      (\{a_i\}_{i\in S_0},e_\nu;\beta_0)\,g^{\nu\mu}\,
   \operatorname{GW}_{A_1,k_1+1}^M
      (e_\mu,\{a_j\}_{j\in S_1};\beta_1),
\end{align*}
这里对 $\nu,\mu$ 采用双重求和约定，并定义
\[
  \varepsilon(S,a)
  :=(-1)^{\#\{j<i\mid i\in S_0,\ j\in S_1,\,
                  \deg(a_i)\deg(a_j)\in2\mathbb Z+1\}}.
\]

\begin{remark}\label{rem:7-5-1}
\begin{enumerate}[label=\textup{(\roman*)},leftmargin=3em]
\item （有效性）公理来自如下事实：该不变量以表示类 $A$ 的 $J$-全纯\BookRecordIndex{T0340}稳定映射之\BookRecordIndex{T0246}模空间
$\overline{\mathcal M}_{0,k}(A;J)$ 为基础。若 $\omega(A)<0$，\BookRecordIndex{T0249}能量恒等式说明这个\BookRecordIndex{T0246}模空间
为空。
\item （对称性）公理来自如下事实：对称群 $S_k$ 通过置换\BookRecordIndex{T0066}标记点，以显然的方式作用在
$\overline{\mathcal M}_{0,k}(A;J)$ 上；\BookRecordIndex{T0283}求值映射
$\operatorname{ev}:\overline{\mathcal M}_{0,k}(A;J)\to M^k$ 与投影
$\pi:\overline{\mathcal M}_{0,k}(A;J)\to\overline{\mathcal M}_{0,k}$ 对此作用等变。
\item （分次）公理来自\BookRecordIndex{T0246}模空间 $\overline{\mathcal M}_{0,k}(A;J)$ 的维数公式，也可以从
（同调）公理\BookRecordIndex{T0326}推出。
\item （同调）公理中的同调类
$\sigma_{A,k}\in H_*(M^k\times\overline{\mathcal M}_{0,k})$，应理解为虚基本类
\[
  [\overline{\mathcal M}_{0,k}(A;J)]^{\mathrm{virt}}
  \in H_*(\overline{\mathcal M}_{0,k}(A;J))
\]
在乘积映射
\[
  \operatorname{ev}_J\times\pi:
  \overline{\mathcal M}_{0,k}(A;J)
  \longrightarrow M^k\times\overline{\mathcal M}_{0,k}
\]
下的上推。在本书的框架内，这条公理仍是启发式陈述，因为只有在很受限的假设下
$\operatorname{ev}_J\times\pi$ 才是\BookRecordIndex{T0337}伪循环；参见习题~\ref{ex:6-7-13} 之前的讨论。正如
那里所说，要在一般情形建立这条公理，需要超出本书范围的分析技术。

不过，如果 $(M,\omega)$ \BookRecordIndex{T0058}半正，那么 $\sigma_{A,k}$ 在投影
$M^k\times\overline{\mathcal M}_{0,k}\to M^k$ 下的上推，在当前框架中有严格意义：它就是
由定理~\ref{thm:7-1-1} 的\BookRecordIndex{T0145}格罗莫夫--威滕伪循环
\[
  \operatorname{ev}_J:\mathcal M^*_{0,k}(A;J)\longrightarrow M^k
\]
所表示的同调类
\[
  \tau_{A,k}\in H_{2n+2c_1(A)+2k-6}(M^k).
\]
再者，若 $T$ 是 $k$-\BookRecordIndex{T0084}带标号树，则 $\sigma_{A,k}$ 与上同调类
$\operatorname{PD}([\overline{\mathcal M}_{0,T}])
\in H^{2e(T)}(\overline{\mathcal M}_{0,k})$ 的配对（或斜积），可以解释为同调类
\[
  \tau_{A,T}\in H_{2n+2c_1(A)+2k-6-2e(T)}(M^k),
\]
它由习题~\ref{ex:6-6-3} 和第~\ref{sec:7-2} 节所讨论的\BookRecordIndex{T0337}伪循环
\[
  \operatorname{ev}_{T,J}:\mathcal M^*_{0,T}(A;J)\longrightarrow M^k
\]
表示。类似地，当 $\#I\ge3$ 时，在\BookRecordIndex{T0066}标记点
$\mathbf w=\{w_i\}_{i\in I}$ 固定的子流形所表示的类
\[
  [Y_{k,I}]\in H_{2(k-\#I)}(\overline{\mathcal M}_{0,k})
\]
的 Poincaré 对偶为
\[
  \operatorname{PD}([Y_{k,I}])
  =\pi_{k,I}^*\operatorname{PD}([\mathrm{pt}])
  \in H^{2\#I-6}(\overline{\mathcal M}_{0,k}),
\]
其中 $\pi_{k,I}:\overline{\mathcal M}_{0,k}\to\overline{\mathcal M}_{0,I}$ 是\BookRecordIndex{T0377}遗忘映射。
所以，假如 $\sigma_{A,k}$ 这样的类存在，它与
$\pi_{k,I}^*\operatorname{PD}([\mathrm{pt}])$ 的配对上推后会得到同调类
\[
  \tau_{A,I}\in H_{2n+2c_1(A)+2(k-\#I)}(M^k),
\]
它由定理~\ref{thm:7-3-1} 的\BookRecordIndex{T0337}伪循环
$\operatorname{ev}_{\mathbf w,J}:\mathcal M_{0,k}(A,\mathbf w;J)\to M^k$ 表示。

注意，我们只在 $\mathbf w$ 属于 $\overline{\mathcal M}_{0,I}$ 的最高\BookRecordIndex{T0072}层时，证明了这个
\BookRecordIndex{T0337}伪循环与 $\mathbf w$ 的选择无关。换言之，由于不能假定 $\sigma_{A,k}$ 存在，我们对
$\tau_{A,I}$ 的定义原则上可能依赖于所选的 $[Y_{k,I}]$ 代表元。下面将看到，（分裂）公理
说明，尽管不能允许完全任意的 $[Y_{k,I}]$ 代表元，至少可以任取一点
$\mathbf w\in\overline{\mathcal M}_{0,I}$。
\item 在\BookRecordIndex{T0058}半正情形，（基本类）、（除子）和（零类）公理很容易从第~\ref{sec:7-1}、
\ref{sec:7-3} 节中\BookRecordIndex{T0143}格罗莫夫--威滕不变量的定义\BookRecordIndex{T0326}推出（见例~\ref{ex:7-1-3}、
\ref{ex:7-3-5} 以及习题~\ref{ex:7-1-5}、\ref{ex:7-1-6}）。下面会为不变量
$\operatorname{GW}$ 和
$\operatorname{GW}^I$ 重述这些公理。
\item \BookRecordIndex{T0143}格罗莫夫--威滕不变量最重要的性质是（分裂）公理。其名称来自如下观察：它可以
解释为与指标集 $\{1,\ldots,k\}$ 的一个分拆相应的同调类 $\sigma_{A,k}$ 分解公式。
公理中的和是有限的，这来自（有效性）和（零类）公理。它的证明用到 \BookRecordIndex{T0003}$J$-全纯曲线的
\BookRecordIndex{T0256}粘合定理；下面以及第~10、11 章将更详细地讨论这一点。现在先注意，它给出的不是计数
全部 $A$-曲线的“通常”\BookRecordIndex{T0143}格罗莫夫--威滕不变量，而是定义域位于
$\overline{\mathcal M}_{0,k}$ 中形如
$\beta=\phi_{S*}(\beta_0\otimes\beta_1)$ 的边界\BookRecordIndex{T0072}层时的受限不变量。最简单的情形是
$k=4$ 且 $\beta_0,\beta_1$ 都等于
$[\mathrm{pt}]=[\overline{\mathcal M}_{0,3}]$。此时
$\phi_{S*}(\beta_0\otimes\beta_1)=[\mathrm{pt}]
\in H_0(\overline{\mathcal M}_{0,4})$，故
\[
  \operatorname{GW}_{A,4}^M
  (a_1,\ldots,a_4;\phi_{S*}(\beta_0\otimes\beta_1))
  =\operatorname{GW}_{A,4}^{M,\{1,2,3,4\}}(a_1,\ldots,a_4),
\]
即\BookRecordIndex{T0185}交比固定的四点不变量。（分裂）公理说明，它化为三点不变量的乘积。此外，如
习题~\ref{ex:7-2-8} 所解释，这个乘积恰是对应于分拆
$\{1,2,3,4\}=S_0\sqcup S_1$ 所给\BookRecordIndex{T0072}层 $T$ 的不变量
$\operatorname{GW}_{A,T}^M(a_1,\ldots,a_4)$。在这个简单情形，\BookRecordIndex{T0072}层 $T$ 是单点，所以也可以
把该乘积看成由\BookRecordIndex{T0337}伪循环 $\operatorname{ev}_{\mathbf w}$ 定义，其中
$\mathbf w$ 不在 $\overline{\mathcal M}_{0,4}$ 的最高\BookRecordIndex{T0072}层。正是在这个意义下，（分裂）公理
与上面 (iv) 中关于（同调）公理的讨论有关。（另请比较习题~\ref{ex:6-7-13} 之前的说明。）
\end{enumerate}
\end{remark}

\begin{remark}\label{rem:7-5-2}
以后有时还需要考虑 $0\le k\le2$ 时带 $k$ 个\BookRecordIndex{T0066}标记点的\BookRecordIndex{T0143}格罗莫夫--威滕不变量
$\operatorname{GW}_{A,k}^M(a_1,\ldots,a_k;\beta)$。由于少于三个\BookRecordIndex{T0066}标记点的亏格零\BookRecordIndex{T0339}稳定曲线
不存在，此时空间 $\overline{\mathcal M}_{0,k}$ 为空，所以这些不变量不属于上述框架。
不过，可以形式地把 $\overline{\mathcal M}_{0,k}$ 看成一个维数为 $2k-6$ 的负维空间，
其基本循环为
$\beta_k\in H_{2k-6}(\overline{\mathcal M}_{0,k})$，并按照注~\ref{rem:7-3-8} 把
$\operatorname{GW}_{A,k}^M(a_1,\ldots,a_k;\beta_k)$ 定义为不变量
$\operatorname{GW}_{A,k}^M(a_1,\ldots,a_k)$；这样（分次）公理依定义成立。此外，当
$k\ge4$ 时，
\[
  \beta_k=\operatorname{PD}
  \bigl(\pi_{0,k}^*\operatorname{PD}(\beta_{k-1})\bigr).
\]
若把这个恒等式形式地延拓到 $k\le3$，那么（除子）公理就化为习题~\ref{ex:7-1-5}
中的陈述，
因而也依定义成立。现在考察（基本类）公理。容易直接验证，当 $A\ne0$ 且 $k\le3$ 时，
\[
  \operatorname{GW}_{A,k}^M(a_1,\ldots,a_{k-1},1;\beta_k)=0.
\]
我们还断言
\[
  \operatorname{GW}_{0,k}^M(a_1,\ldots,a_k;\beta_k)=0,
  \qquad k=0,1,2.
\]
这来自维数理由。这个不变量应当计数与 $a_1,\ldots,a_k$ 的 Poincaré 对偶循环相交的
常值球面。（分次）公理蕴含
\[
  \sum_i\deg(a_i)=\dim M+2k-6<\dim M.
\]
因此，这些循环不相交，不变量为零。（比较例~\ref{ex:7-1-3}，那里从另一观点得到了类似结论。）
这同时延拓了（基本类）与（零类）公理。（有效性）和（对称性）公理仍然成立，而
（同调）和（分裂）公理不适用。下面的习题说明，这些额外的不变量中有些并不为零。
\end{remark}

\begin{exercise}\label{ex:7-5-3}
用 $L$ 表示 $\mathbb{CP}^n$ 中由一条直线表示的类，并令
$c\in H^2(\mathbb{CP}^n)$ 满足 $c(L)=1$。验证：对每个 $n\ge1$，
\[
  \operatorname{GW}_{L,0}^{\mathbb{CP}^1}(\beta_0)
  =\operatorname{GW}_{L,1}^{\mathbb{CP}^1}(c;\beta_1)=1,
  \qquad
  \operatorname{GW}_{L,2}^{\mathbb{CP}^n}(c^n,c^n;\beta_2)=1,
\]
并且当 $k=0,1$、$n\ge2$ 时
$\operatorname{GW}_{L,k}^{\mathbb{CP}^n}=0$。
\end{exercise}

\paragraph{作为积分的\BookRecordIndex{T0143}格罗莫夫--威滕不变量。}
在\BookRecordIndex{T0058}半正情形，只要用微分形式表示各类 $a_i$，而这些微分形式支撑在其对偶循环的足够小
邻域中，并且这些对偶循环与\BookRecordIndex{T0246}模空间 $\mathcal M^*_{0,k}(A;J)$ 上的\BookRecordIndex{T0283}求值映射横截，就能严格
实现\BookRecordIndex{T0143}格罗莫夫--威滕不变量的积分表述。此时，乘积 $a_1\wedge\cdots\wedge a_k$ 在\BookRecordIndex{T0283}求值映射
下的拉回具有紧支撑，故积分有意义。为了证明积分与形式的选择无关，需要证明任意\BookRecordIndex{T0403}正合形式
的拉回积分为零，即
\[
  \int_{\mathcal M^*_{0,k}(A;J)}\operatorname{ev}^*d\sigma=0.
\]
直观地说，这应当成立，因为 $\mathcal M^*_{0,k}(A;J)$ 的边界余维为二，所以 $\sigma$ 在
边界上的积分应当为零。然而，要严格说明这一点并不容易。在相交理论的做法中，可以借助
\BookRecordIndex{T0296}萨德定理避开这个问题。

积分表述突出了\BookRecordIndex{T0143}格罗莫夫--威滕不变量与其他规范理论不变量的相似之处，例如可参见
Ruan \srcite{342}、Givental \srcite{151} 和 Manin \srcite{286}。

\paragraph{代数重述。}
Kontsevich 和 Manin 在 \srcite{216} 中把\BookRecordIndex{T0143}格罗莫夫--威滕不变量描述为同态
\[
  I_{A,k}:H^*(M)^{\otimes k}\longrightarrow
  H^*(\overline{\mathcal M}_{0,k}),
  \qquad k\ge2.
\]
把同态 $I_{A,k}$ 与 $H_*(\overline{\mathcal M}_{0,k})$ 中的同调类配对，就得到不变量
$\operatorname{GW}_{A,k}$：
\[
  \operatorname{GW}_{A,k}(a_1,\ldots,a_k;\beta)
  =\bigl\langle I_{A,k}(a_1,\ldots,a_k),\beta\bigr\rangle.
\]
于是，\BookRecordIndex{T0254}粘合规则可以看成赋予 $H^*(M)$ 一个\BookRecordIndex{T0299}上同调场论的结构（见
\srcite{286}，定义 0.3.1）。第二种重述强调
$\overline{\mathcal M}_{0,k+1}$ 的同调在 $M$ 的上同调上的作用；它来自对同态 $I_{A,k}$
作部分对偶。所得多重线性映射
\[
  m_{A,k}:H_*(\overline{\mathcal M}_{0,k+1})
  \otimes H^*(M)^{\otimes k}\longrightarrow H^*(M)
\]
由下式定义：
\[
  \bigl\langle m_{A,k}(\beta;a_1,\ldots,a_k),
                 \operatorname{PD}(a_{k+1})\bigr\rangle
  :=\bigl\langle\beta,I_{A,k+1}(a_1,\ldots,a_{k+1})\bigr\rangle.
\]
特别地，若 $k=2$，则 $\overline{\mathcal M}_{0,3}$ 是一点；取
$\beta=[\overline{\mathcal M}_{0,3}]$，得到双线性映射
\[
  m_A:H^*(M)\otimes H^*(M)\longrightarrow H^*(M).
\]
（分裂）公理的一个推论是：把这些 $m_A$ 对 $A$ 求和，会在 $H^*(M)$ 上定义结合乘法，
常称为\BookRecordIndex{T0349}小量子乘法。\footnote{为了严格定义这个乘法，必须用 Novikov 环扩张
$H=H^*(M)$ 的系数，以记录 $A$；见第~11 章。}物理学家猜想了这一乘法的结合性；事实上，
这只是他们认为各映射 $m_{A,k}$ 应满足的一族方程中的第一个。最容易的情形
（$\beta=[\overline{\mathcal M}_{0,3}]$ 且 $(M,\omega)$ \BookRecordIndex{T0058}半正）最先由 Ruan--Tian 在
\srcite{346} 中证明；其他类 $\beta$ 及高亏格情形在 \srcite{348} 中处理。不过，一般
$(M,\omega)$ 的情形必须等到\BookRecordIndex{T0370}虚模循环得到理解之后才能解决。第~11 章会更详细地描述
量子乘法，第~10 章会证明其结合性（在\BookRecordIndex{T0058}半正情形）。

\begin{remark}\label{rem:7-5-4}
$\overline{\mathcal M}_{0,k}$ 上有一些与\BookRecordIndex{T0143}格罗莫夫--威滕不变量发生有趣相互作用的自然上同调
类，它们是典范线丛
\[
  L_i\longrightarrow\overline{\mathcal M}_{0,k},
  \qquad i=1,\ldots,k
\]
的第一 Chern 类。$L_i$ 在 $\mathbf w$ 处的纤维，是 $\mathbf w$ 对应的\BookRecordIndex{T0339}稳定曲线在第 $i$ 个
\BookRecordIndex{T0066}标记点处的余切空间。等价地，把 $\overline{\mathcal M}_{0,k}$ 看作
$\overline{\mathcal M}_{0,k+1}$ 中由第 $i$ 个\BookRecordIndex{T0066}标记点决定的子流形时，$L_i$ 是它的余\BookRecordIndex{T0117}法丛。
$L_i$ 的第一 Chern 类给出所谓的\BookRecordIndex{T0379}引力祖先；\BookRecordIndex{T0340}稳定映射\BookRecordIndex{T0246}模空间上类似线丛则给出\BookRecordIndex{T0378}引力后裔。
更多信息见 Coates \srcite{72}、Givental \srcite{151} 和 Manin \srcite{286}。
\end{remark}

\paragraph{$\operatorname{GW}$ 与 $\operatorname{GW}^I$ 的公理。}
现在更仔细地考察这些公理，并为第~\ref{sec:7-1}、\ref{sec:7-3} 节定义的不变量
$\operatorname{GW}$ 和 $\operatorname{GW}^I$ 重述它们。如果能证明具有上述全部性质的
\BookRecordIndex{T0143}格罗莫夫--威滕不变量确实存在，那么这些 $\operatorname{GW}$ 和
$\operatorname{GW}^I$ 应有如下解释：
\[
  \operatorname{GW}_{A,k}^M(a_1,\ldots,a_k)
  =:\operatorname{GW}_{A,k}^M
     (a_1,\ldots,a_k;[\overline{\mathcal M}_{0,k}]),
\]
而对满足 $\#I\ge3$ 的指标集 $I\subset\{1,\ldots,k\}$，
\[
  \operatorname{GW}_{A,k}^{M,I}(a_1,\ldots,a_k)
  =:\operatorname{GW}_{A,k}^M(a_1,\ldots,a_k;\beta_{k,I}).
\]
这里的同调类
$\beta_{k,I}\in H_*(\overline{\mathcal M}_{0,k})$ 定义为
\[
  \beta_{k,I}
  :=\operatorname{PD}\bigl(
    \pi_{k,I}^*\operatorname{PD}([\mathrm{pt}])\bigr)
  \in H_{2(k-\#I)}(\overline{\mathcal M}_{0,k}),
\]
其中
\[
  \pi_{k,I}:\overline{\mathcal M}_{0,k}
  \longrightarrow\overline{\mathcal M}_{0,I}
\]
是遗忘不属于指标集 $I$ 的\BookRecordIndex{T0066}标记点的映射。这意味着，$\beta_{k,I}$ 可以由\BookRecordIndex{T0150}固定标记点
$\mathbf w:=\{w_i\}_{i\in I}$ 所得的 $\overline{\mathcal M}_{0,k}$ 中的循环表示。事实上，
下面的引理~\ref{lem:7-5-5} 将证明，它可由 $\pi_{k,I}$ 的任意\BookRecordIndex{T0409}正则值
$\mathbf w\in\overline{\mathcal M}_{0,I}$ 的逆像 $\pi_{k,I}^{-1}(\mathbf w)$ 表示。注意，
当 $\#I=3$ 时，$\beta_{k,I}$ 是 $\overline{\mathcal M}_{0,k}$ 的基本循环。

现在说明 $\operatorname{GW}$ 和 $\operatorname{GW}^I$ 确实具有预期性质。证明以同调类
$\beta_{k,I}$ 的性质为基础。通常如第~\ref{sec:7-3} 节那样，选取
$\mathbf w\in\overline{\mathcal M}_{0,I}$ 的开\BookRecordIndex{T0072}层
$\mathcal M_{0,I}$，也就是令 $w_i$（$i\in I$）为 $S^2$ 上两两不同的点。但后面的论证中，
选取一个更典范、各 $w_i$ 尽可能分开的 $\mathbf w$ 会更方便；换言之，我们让
$\mathbf w$ 以具有尽可能多分支的树为模型。为此，写
\[
  I=:\{i_0,\ldots,i_\ell\},
  \qquad \ell:=\#I-1,
\]
并令 $T$ 为具有 $\ell-1$ 个顶点
$\alpha_1,\ldots,\alpha_{\ell-1}$ 的树，其中
$\alpha_jE\alpha_{j+1}$（$j=1,\ldots,\ell-2$）。定义
$\mathbf w_I\in\overline{\mathcal M}_{0,I}$ 为对应于 $I$-\BookRecordIndex{T0084}带标号树 $(T,\Lambda,E)$ 的唯一
元素，其中
\begin{equation}\label{eq:7-5-1}
  i_0,i_1\in\Lambda_{\alpha_1},\qquad
  i_j\in\Lambda_{\alpha_j},\qquad
  i_{\ell-1},i_\ell\in\Lambda_{\alpha_{\ell-1}}.
\end{equation}
因此，对 $0\le j_0<j_1<j_2<j_3\le\ell$，$\mathbf w_I$ 的分量满足
\[
  w_{i_{j_0}i_{j_1}i_{j_2}i_{j_3}}=\infty.
\]

\begin{lemma}\label{lem:7-5-5}
设 $I\subset\{1,\ldots,k\}$ 且 $\#I\ge3$。
\begin{enumerate}[label=\textup{(\roman*)},leftmargin=3em]
\item 类 $\beta_{k,I}$ 可以由循环
$Y_{k,I}:=\pi_{k,I}^{-1}(\mathbf w_I)$ 表示。
\item 用 $\pi_{0,k}:\overline{\mathcal M}_{0,k}\to
\overline{\mathcal M}_{0,k-1}$ 表示遗忘最后一个\BookRecordIndex{T0066}标记点的投影。当 $k\notin I$ 时，
\[
  \beta_{k,I}
  =\operatorname{PD}\bigl(\pi_{0,k}^*\operatorname{PD}(\beta_{k-1,I})\bigr).
\]
此外，
\[
  (\pi_{0,k})_*\beta_{k,I}
  =\begin{cases}
    \beta_{k-1,I\setminus\{k\}},&k\in I,\ \#I\ge4,\\
    0,&k\notin I.
  \end{cases}
\]
\end{enumerate}
\end{lemma}

\begin{proof}
设 $f:X\to Y$ 是闭定向流形之间的光滑映射。利用 Poincaré 对偶，可以定义\BookRecordIndex{T0323}同调拉回
（或“上惊叹”）
\[
  f^!:H_*(Y)\longrightarrow H_*(X),
  \qquad f^!(a):=\operatorname{PD}_Xf^*\operatorname{PD}_Y(a).
\]
几何上，如果 $H_*(Y)$ 中的一个同调类由与 $f$ 横截的定向子流形 $Z\subset Y$ 表示，
则它在 $f^!$ 下的像由定向子流形 $f^{-1}(Z)\subset X$ 表示。\BookRecordIndex{T0323}同调拉回具有函子性：
$(f\circ g)^!=g^!\circ f^!$。用这种语言，
$\beta_{k,I}=\pi_{k,I}^!([\mathrm{pt}])$。

为证明 (i)，需要验证 $\mathbf w_I$ 是 $\pi_{k,I}$ 的\BookRecordIndex{T0409}正则值。只须证明：对每个
$\mathbf z\in Y_{k,I}$，$\overline{\mathcal M}_{0,I}$ 中 $\mathbf w_I$ 附近的局部坐标都能
扩充为 $\overline{\mathcal M}_{0,k}$ 中 $\mathbf z$ 附近的一组局部坐标。每个这样的点
$\mathbf z$ 都属于某个子流形 $\overline{\mathcal M}_{0,T,\Lambda}$，其中 $\Lambda$ 是满足
式~\eqref{eq:7-5-1} 的 $T$ 的一个 $k$-标号。因此，断言来自定理 D.4.2 的证明。
$\mathbf w_I$ 附近 $\overline{\mathcal M}_{0,I}$ 上的坐标对应于指标集 $I$ 的商树中的边；
再按照定理 D.4.2 证明中的方法选取 $\overline{\mathcal M}_{0,T,\Lambda}$ 上的坐标，就能把
它们扩充为 $\mathbf z$ 附近 $\overline{\mathcal M}_{0,k}$ 上的一组坐标。这就证明了 (i)。

为证明 (ii)，先设 $k\notin I$。此时
$\pi_{k,I}=\pi_{k-1,I}\circ\pi_{0,k}$，故关于 $\beta_{k,I}$ 的恒等式来自\BookRecordIndex{T0323}同调拉回的函子性。
再考察上推 $(\pi_{0,k})_*\beta_{k,I}$。显然
$\pi_{0,k}(Y_{k,I})=Y_{k-1,I\setminus\{k\}}$，其中
$Y_{k,I}:=\pi_{k,I}^{-1}(\mathbf w_I)$。若 $k\in I$ 且 $\#I\ge4$，则两个循环具有相同维数
$2(k-\#I)$，并且 $\pi_{0,k}$ 在 $Y_{k,I}$ 上的限制是从 $Y_{k,I}$ 到
$Y_{k-1,I\setminus\{k\}}$ 的全纯微分同胚。\sourcefootnote{原书把这里两个循环的维数印作
$2(k-2\#I)$；但由 $Y_{k,I}=\pi_{k,I}^{-1}(\mathbf w_I)$ 的余维以及本节前面对
$\beta_{k,I}$ 次数的定义，维数均为 $2(k-\#I)$，故据此补正。}这证明了第二个断言。最后，
若 $k\notin I$，则 $\pi_{0,k}$ 把 $Y_{k,I}$ 映到一个较低维流形，因而最后一个断言成立。
\end{proof}

下面是对不变量 $\operatorname{GW}^M$ 的（除子）、（基本类）与（零类）公理的重述。

\begin{proposition}\label{prop:7-5-6}
设 $(M,\omega)$ 是紧\BookRecordIndex{T0058}半正\BookRecordIndex{T0358}辛流形，$A\in H_2(M;\mathbb Z)$，$k\ge1$，并且
$a_1,\ldots,a_k\in H^*(M)$。则有：
\begin{description}[leftmargin=2em,style=nextline]
\item[（除子）] 若 $(A,k)\ne(0,3)$ 且 $\deg(a_k)=2$，则
\[
  \operatorname{GW}_{A,k}^M(a_1,\ldots,a_k)
  =\operatorname{GW}_{A,k-1}^M(a_1,\ldots,a_{k-1})\int_Aa_k.
\]
\item[（基本类）] 若 $(A,k)\ne(0,3)$，则
\[
  \operatorname{GW}_{A,k}^M(a_1,\ldots,a_{k-1},1)=0.
\]
\item[（零类）] 若 $k\ne3$，则 $\operatorname{GW}_{0,k}^M=0$；若 $k=3$，则
\[
  \operatorname{GW}_{0,3}^M(a_1,a_2,a_3)
  =\int_Ma_1\smile a_2\smile a_3.
\]
\end{description}
\end{proposition}

\begin{proof}
留作习题。（零类）公理的详细讨论见例~\ref{ex:7-1-3}、\ref{ex:7-3-5}；另两条公理背后的
几何直观见习题~\ref{ex:7-1-6}、\ref{ex:7-1-5}。
\end{proof}

下面给出 $\operatorname{GW}^{M,I}$ 的相应结果。注意，现在基本类公理更为微妙，这反映了
完整不变量族 $\operatorname{GW}_{A,k}^M(a_1,\ldots,a_k;\beta)$ 的公理结构。

\begin{proposition}\label{prop:7-5-7}
设 $(M,\omega)$ 是紧\BookRecordIndex{T0058}半正\BookRecordIndex{T0358}辛流形，$A\in H_2(M;\mathbb Z)$，$k\ge3$，
$I\subset\{1,\ldots,k\}$ 且 $\#I\ge3$，并且 $a_1,\ldots,a_k\in H^*(M)$。则有：
\begin{description}[leftmargin=2em,style=nextline]
\item[（除子）] 若 $k\notin I$ 且 $\deg(a_k)=2$，则
\[
  \operatorname{GW}_{A,k}^{M,I}(a_1,\ldots,a_k)
  =\operatorname{GW}_{A,k-1}^{M,I}(a_1,\ldots,a_{k-1})\int_Aa_k.
\]
\item[（基本类）] 若 $k\in I$ 且 $\#I\ge4$，则
\[
  \operatorname{GW}_{A,k}^{M,I}(a_1,\ldots,a_{k-1},1)
  =\operatorname{GW}_{A,k-1}^{M,I\setminus\{k\}}(a_1,\ldots,a_{k-1}).
\]
若 $k\notin I$，则
$\operatorname{GW}_{A,k}^{M,I}(a_1,\ldots,a_{k-1},1)=0$。
\item[（零类）] 若 $\#I<k$，则 $\operatorname{GW}_{0,k}^{M,I}=0$；若 $\#I=k$，则
\[
  \operatorname{GW}_{0,k}^{M,\{1,\ldots,k\}}(a_1,\ldots,a_k)
  =\int_Ma_1\smile\cdots\smile a_k.
\]
\end{description}
\end{proposition}

\begin{proof}
留作习题。
\end{proof}

\begin{exercise}\label{ex:7-5-8}
为不变量 $\operatorname{GW}$ 和 $\operatorname{GW}^I$ 表述并证明（对称性）公理。由于置换群
$S_k$ 通过同痕于恒等映射的微分同胚作用在 $\overline{\mathcal M}_{0,k}(A;J)$ 上，公式中的
符号来自 $S_k$ 在乘积 $M^k$ 的上同调上的作用。
\end{exercise}

下面考察（分裂）公理中出现的乘积映射
\[
  \phi_S:\overline{\mathcal M}_{0,k_0+1}
  \times\overline{\mathcal M}_{0,k_1+1}
  \longrightarrow\overline{\mathcal M}_{0,k}.
\]
直观地说，它把 $\overline{\mathcal M}_{0,k_0+1}$ 中\BookRecordIndex{T0339}稳定曲线的最后一个\BookRecordIndex{T0066}标记点与
$\overline{\mathcal M}_{0,k_1+1}$ 中\BookRecordIndex{T0339}稳定曲线的第一个\BookRecordIndex{T0066}标记点等同，从而把两条\BookRecordIndex{T0339}稳定曲线连接
起来。所得曲线在连接处有一个\BookRecordIndex{T0191}节点，因而只剩 $k_0+k_1$ 个\BookRecordIndex{T0066}标记点。现在说明：当某些\BookRecordIndex{T0066}标记点
固定时，必须假定将要\BookRecordIndex{T0254}粘合的两个点也固定；否则，第一条曲线上的固定点与第二条曲线上的
固定点之间的关系可能变化。这正是下文两个子集 $I_{S_0},I_{S_1}$ 的定义所表达的内容。

设 $I_0,I_1\subset\{1,\ldots,k\}$ 是不交子集，满足
\[
  \#I_0\ge2,\qquad \#I_1\ge2,
\]
并记
\[
  I:=I_0\cup I_1.
\]
用 $\mathcal S(I_0,I_1)$ 表示指标集 $\{1,\ldots,k\}$ 的所有分拆
$S=(S_0,S_1)$ 所成的集合，其中 $I_0\subset S_0$ 且 $I_1\subset S_1$。给定
\[
  S=(S_0,S_1)\in\mathcal S(I_0,I_1),
\]
令 $k_i:=\#S_i$，并令
\[
  \sigma_0:\{1,\ldots,k_0\}\longrightarrow\{1,\ldots,k\},
  \qquad
  \sigma_1:\{2,\ldots,k_1+1\}\longrightarrow\{1,\ldots,k\}
\]
为满足 $\operatorname{im}\sigma_0=S_0$、$\operatorname{im}\sigma_1=S_1$ 的唯一保序单射。
定义 $I_{S_i}\subset\{1,\ldots,k_i+1\}$ 为
\[
  I_{S_0}:=\sigma_0^{-1}(I)\cup\{k_0+1\},
  \qquad
  I_{S_1}:=\{1\}\cup\sigma_1^{-1}(I).
\]

\begin{lemma}[Keel]\label{lem:7-5-9}
\[
  \beta_{k,I}
  =\sum_{S\in\mathcal S(I_0,I_1)}
    \phi_{S*}\bigl(
      \beta_{k_0+1,I_{S_0}}\otimes\beta_{k_1+1,I_{S_1}}
    \bigr).
\]
\end{lemma}

\begin{proof}
先明确描述映射 $\phi_S$。它由
\[
  \phi_S(\mathbf w^0,\mathbf w^1):=\mathbf w^S
\]
给出，其中 $\mathbf w^S$ 由下列条件定义：当
$i,i'\in S_0$ 且 $j,j'\in S_1$ 时，
\[
  w^S_{ii'jj'}:=\infty;
\]
当 $i,i',i''\in\{1,\ldots,k_0\}$ 且
$j,j',j''\in\{2,\ldots,k_1+1\}$ 时，
\[
  w^S_{\sigma_0(i)\sigma_0(i')\sigma_0(i'')\sigma_1(j)}
  :=w^0_{ii'i''\,k_0+1},
  \qquad
  w^S_{\sigma_0(i)\sigma_1(j)\sigma_1(j')\sigma_1(j'')}
  :=w^1_{1jj'j''};
\]
并且，当 $i_0,i_1,i_2,i_3\in S_i$、$i=0,1$ 时，
\[
  w^S_{i_0i_1i_2i_3}
  :=w^i_{\sigma_i^{-1}(i_0)\sigma_i^{-1}(i_1)
              \sigma_i^{-1}(i_2)\sigma_i^{-1}(i_3)}.
\]

现在考察引理~\ref{lem:7-5-5} 证明中定义的子流形
$Y_{k,I}\subset\overline{\mathcal M}_{0,k}$。我们断言
\[
  Y_{k,I}
  =\bigcup_{S\in\mathcal S(I_0,I_1)}
    \phi_S\bigl(Y_{k_0+1,I_{S_0}}\times Y_{k_1+1,I_{S_1}}\bigr).
\]
从前面给出的\BookRecordIndex{T0257}粘合映射 $\phi_S$ 的几何描述来看，这是显然的。若要从上面的逐坐标描述验证，
只须注意
\[
  Y_{k,I}:=\bigl\{\mathbf w\in\overline{\mathcal M}_{0,k}\bigm|
  i_0,i_1,i_2,i_3\in I,\ i_0<i_1<i_2<i_3
  \Longrightarrow w_{i_0i_1i_2i_3}=\infty\bigr\}.
\]
子流形 $Y_{k,I}$ 与
$Y_{k_0+1,I_{S_0}}\times Y_{k_1+1,I_{S_1}}$ 具有相同维数 $2(k-\#I)$。此外，
$\phi_S$ 在后一子流形上的限制，是到 $Y_{k,I}$ 的若干连通分支之并上的全纯微分同胚。
这就证明了引理。
\end{proof}

把引理~\ref{lem:7-5-9} 与（分裂）公理合起来，就得到\BookRecordIndex{T0143}格罗莫夫--威滕不变量
$\operatorname{GW}^{M,I}$ 的如下\BookRecordIndex{T0256}粘合定理。

\begin{theorem}[\BookRecordIndex{T0254}粘合]\label{thm:7-5-10}
设 $(M,\omega)$ \BookRecordIndex{T0058}半正。则
\begin{align*}
  \operatorname{GW}_{A,k}^{M,I}(a_1,\ldots,a_k)
  &={}
  \sum_{S\in\mathcal S(I_0,I_1)}\varepsilon(S,a)
  \sum_{A_0+A_1=A}\sum_{\nu,\mu}
  \operatorname{GW}_{A_0,k_0+1}^{M,I_{S_0}}
     (\{a_i\}_{i\in S_0},e_\nu)\,g^{\nu\mu}\\
  &\hspace{7em}\cdot
  \operatorname{GW}_{A_1,k_1+1}^{M,I_{S_1}}
     (e_\mu,\{a_j\}_{j\in S_1}).
\end{align*}
特别地，
\begin{align*}
  &\operatorname{GW}_{A,k}^{M,\{1,2,3,4\}}(a_1,\ldots,a_k)\\
  &\quad={}
  \sum_{S\in\mathcal S(\{1,2\},\{3,4\})}\varepsilon(S,a)
  \sum_{A_0+A_1=A}\sum_{\nu,\mu}
  \operatorname{GW}_{A_0,k_0+1}^M(\{a_i\}_{i\in S_0},e_\nu)
  \,g^{\nu\mu}\\
  &\qquad\cdot
  \operatorname{GW}_{A_1,k_1+1}^M(e_\mu,\{a_j\}_{j\in S_1}),
\end{align*}
而当 $k_0+k_1=k$ 且 $k_0,k_1\ge2$ 时，
\begin{align*}
  &\operatorname{GW}_{A,k}^{M,\{1,\ldots,k\}}(a_1,\ldots,a_k)\\
  &\quad={}
  \sum_{A_0+A_1=A}\sum_{\nu,\mu}
  \operatorname{GW}_{A_0,k_0+1}^{M,\{1,\ldots,k_0+1\}}
     (a_1,\ldots,a_{k_0},e_\nu)\,g^{\nu\mu}\,
  \operatorname{GW}_{A_1,k_1+1}^{M,\{1,\ldots,k_1+1\}}
     (e_\mu,a_{k_0+1},\ldots,a_k).
\end{align*}
\end{theorem}

为说明这些规则的含义，考察
\[
  k=\#I=4
\]
的情形。此时 $\overline{\mathcal M}_{0,4}$ 是二维球面，其中有三个\BookRecordIndex{T0320}特殊点
$\mathbf w_0,\mathbf w_1,\mathbf w_\infty$，对应于三棵有两个顶点的树。若
$\mathbf w$ 不等于这三个\BookRecordIndex{T0320}特殊点中的任何一个，则
$\operatorname{GW}_{A,4}^{M,\{1,2,3,4\}}(a_1,\ldots,a_4)$ 计数带四个\BookRecordIndex{T0150}固定标记点
$w_i$ 的 $A$-曲线，其中各 $w_i$ 通过与 $a_i$ 对偶的循环。让其中两个\BookRecordIndex{T0066}标记点彼此靠拢，
等价地说，考察 $\overline{\mathcal M}_{0,4}$ 的三个\BookRecordIndex{T0320}特殊点之一，就可以说明同一个不变量
计数具有两个分支的\BookRecordIndex{T0340}稳定映射，每个分支通过两个循环。要证明这两种描述确实给出同一不变量，
必须证明一个\BookRecordIndex{T0256}粘合定理，说明怎样由一对分支构造一条 $A$-曲线。这个结果将在第~10 章证明。

\paragraph{Kontsevich 递推公式。}
现在说明定理~\ref{thm:7-5-10} 的重要性：用它计算 $\mathbb{CP}^2$ 中通过
$3d-1$ 个一般点的 $d$ 次\BookRecordIndex{T0386}有理曲线数
\[
  N_d:=\operatorname{GW}_{dL,3d-1}^{\mathbb{CP}^2}
       (c^2,\ldots,c^2).
\]

\begin{proposition}[Kontsevich]\label{prop:7-5-11}
各数 $N_d$ 由初始条件 $N_1=1$ 和递推公式
\begin{equation}\label{eq:7-5-2}
  N_d=\sum_{k+\ell=d}N_kN_\ell
  \left(
    k^2\ell^2\binom{3d-4}{3k-2}
    -k^3\ell\binom{3d-4}{3k-1}
  \right),
  \qquad d\ge2,
\end{equation}
决定。最初几个值为
\[
  N_2=1,\qquad N_3=12,\qquad N_4=620,\qquad N_5=87{,}304,
\]
\[
  N_6=26{,}312{,}976,\qquad N_7=14{,}616{,}808{,}192.
\]
\end{proposition}

\begin{proof}
由于 $\mathbb{CP}^2$ 的上同调均为偶数次，定理~\ref{thm:7-5-10} 中所有符号都是正号。
考察 $H^*(\mathbb{CP}^2)$ 的标准基
\[
  e_\nu:=c^\nu,\qquad \nu=0,1,2.
\]
于是
\[
  g^{20}=g^{02}=g^{11}=1,
\]
而对 $\nu,\mu$ 的所有其他取值，$g^{\nu\mu}=0$。所以\BookRecordIndex{T0254}粘合公式中有贡献的三对
$(e_\nu,e_\mu)$ 是 $(1,c^2)$、$(c,c)$ 和 $(c^2,1)$。我们将取
$I=\{1,2,3,4\}$ 应用\BookRecordIndex{T0254}粘合公式，使 $I_0,I_1$ 各有两个元素；因而其右端的不变量都是
$\operatorname{GW}^I$ 的形式，其中 $\#I=3$。因此，由（基本类）公理，
任何含 $1$ 的项都为零，除非相应的类 $A_i$ 为零。

下一个观察是，由命题~\ref{prop:7-5-6}，
\begin{equation}\label{eq:7-5-3}
  \operatorname{GW}_{dL,3d+k}^{\mathbb{CP}^2}
  (\underbrace{c,\ldots,c}_{k+1},c^2,\ldots,c^2)
  =d^{k+1}N_d,
\end{equation}
这里前 $k+1$ 个自变量等于 $c$。若 $c$ 型自变量的数目不同，则由维数理由，该不变量为零。

现在计数带四个\BookRecordIndex{T0150}固定标记点、通过两条直线 $\ell_1,\ell_2$ 和
$3d-2$ 个点 $p_1,\ldots,p_{3d-2}$ 的曲线。\BookRecordIndex{T0254}粘合规则说明，当把指标集
$\{1,2,3,4\}$ 分成两对 $I_0:=\{1,2\}$ 与 $I_1:=\{3,4\}$ 时，这个计数怎样分解。
约束 $\operatorname{PD}(a_i)$ 本质上有两种不同的分配方式：或者要求
$I_0$ 中的\BookRecordIndex{T0066}标记点 $w_1,w_2$ 都位于直线上，或者要求其中一个位于直线上、另一个位于一点上。
如下计算表明，这就给出了递推公式。

应用定理~\ref{thm:7-5-10}，取指标集
\[
  I_0:=\{1,2\},\qquad I_1:=\{3,4\},
\]
并把约束排序，使 $w_1,w_2$ 位于直线上，得到
\begin{align*}
 &\operatorname{GW}_{dL,3d}^{\mathbb{CP}^2,\{1,2,3,4\}}
   (c,c,c^2,\ldots,c^2)\\
 &=\sum_{d_0+d_1=d}\binom{3d-4}{3d_0}
   \operatorname{GW}_{d_0L,3d_0+3}^{\mathbb{CP}^2}
     (1,c,c,c^2,\ldots,c^2)
   \operatorname{GW}_{d_1L,3d_1-1}^{\mathbb{CP}^2}
     (c^2,\ldots,c^2)\\
 &\quad+
   \sum_{d_0+d_1=d}\binom{3d-4}{3d_0-1}
   \operatorname{GW}_{d_0L,3d_0+2}^{\mathbb{CP}^2}
     (c,c,c,c^2,\ldots,c^2)
   \operatorname{GW}_{d_1L,3d_1}^{\mathbb{CP}^2}
     (c,c^2,\ldots,c^2)\\
 &=N_d+
   \sum_{d_0+d_1=d}\binom{3d-4}{3d_0-1}
     d_0^3d_1N_{d_0}N_{d_1}.
\end{align*}
第一个等号在\BookRecordIndex{T0254}粘合公式中展开代表对角线的三对
$(e_\nu,e_\mu)=(1,c^2),(c,c),(c^2,1)$，并使用了式~\eqref{eq:7-5-3}。第一个二项式系数
计数指标集 $\{1,\ldots,3d\}$ 的分拆，其中
$1,2\in S_0$、$3,4\in S_1$，且
\[
  \#S_0=3d_0+2,\qquad \#S_1=3d_1-2.
\]
（若这些数值条件不成立，由式~\eqref{eq:7-5-3}，不变量为零。）第二个二项式系数是
\[
  \#S_0=3d_0+1,\qquad \#S_1=3d_1-1
\]
时相应的数目。$(\nu,\mu)=(c^2,1)$ 的第三项由命题~\ref{prop:7-5-6} 的（基本类）公理
为零，这解释了第一个等号。第二个等号也使用了同一事实：由于
\[
  \operatorname{GW}_{d_0L,3d_0+3}^{\mathbb{CP}^2}
  (1,c,c,c^2,\ldots,c^2)=0
\]
除非 $d_0=0$，第一重和唯一的贡献计数如下\BookRecordIndex{T0340}稳定映射：它的第一个分支是恰好带两个\BookRecordIndex{T0066}标记点
$w_1,w_2$、映到两条直线的交点 $\ell_1\cap\ell_2$ 的\BookRecordIndex{T0382}幽灵分支；第二个分支是一条
$d_1=d$ 次曲线，通过 $3d-1$ 个点
$\ell_1\cap\ell_2,p_1,\ldots,p_{3d-2}$。第二重和的求值使用了（除子）公理。

如上所述，置换前四个 $a_i$ 不会改变不变量，但会使约束在分裂中有不同的分配，因为现在
$w_1,w_3$ 位于两条直线上。与上面相同的论证给出
\begin{align*}
  &N_d+
   \sum_{d_0+d_1=d}\binom{3d-4}{3d_0-1}
      d_0^3d_1N_{d_0}N_{d_1}\\
  &=\operatorname{GW}_{dL,3d}^{\mathbb{CP}^2,\{1,2,3,4\}}
      (c,c^2,c,c^2,\ldots,c^2)\\
  &=\sum_{d_0+d_1=d}\binom{3d-4}{3d_0-2}
    \operatorname{GW}_{d_0L,3d_0+1}^{\mathbb{CP}^2}
      (c,c,c^2,\ldots,c^2)
    \operatorname{GW}_{d_1L,3d_1+1}^{\mathbb{CP}^2}
      (c,c,c^2,\ldots,c^2)\\
  &=\sum_{d_0+d_1=d}\binom{3d-4}{3d_0-2}
      d_0^2d_1^2N_{d_0}N_{d_1}.
\end{align*}
这就证明了命题~\ref{prop:7-5-11} 的递推公式。
\end{proof}

\begin{exercise}\label{ex:7-5-12}
用式~\eqref{eq:7-5-2} 验证 $N_2=1$、$N_3=12$。在两种情形都明确描述计算中出现的\BookRecordIndex{T0340}稳定映射
的几何形态。注意，同一种几何构型可能贡献多次：一方面，它可能以不同方式与约束相互作用；
另一方面，它可能有不同的参数化。例如，一条直线与一条圆锥曲线的并可按两种方式参数化为
三次\BookRecordIndex{T0340}稳定映射，因为定义域中的唯一\BookRecordIndex{T0191}节点可以映到直线与圆锥曲线的两个交点中的任一个。
\end{exercise}

Kontsevich 和 Manin 在 \srcite{216} 中还说明，原则上怎样由\BookRecordIndex{T0256}粘合定理计算高维射影空间的
\BookRecordIndex{T0143}格罗莫夫--威滕不变量。他们引入数
\[
  N_d(k_2,\ldots,k_n)
  :=\operatorname{GW}_{dL,k}^{\mathbb{CP}^n}
  (\underbrace{c^2,\ldots,c^2}_{k_2},
   \underbrace{c^3,\ldots,c^3}_{k_3},\ldots,
   \underbrace{c^n,\ldots,c^n}_{k_n}),
\]
其中 $k:=k_2+\cdots+k_n$，上同调类 $c^\nu$ 出现 $k_\nu$ 次。例如，
$N_d(k_2,k_3)$ 是 $\mathbb{CP}^3$ 中通过 $k_2$ 条直线和 $k_3$ 个点的 $d$ 次\BookRecordIndex{T0386}有理曲线数。
只有当这些上同调类满足维数条件~\eqref{eq:7-1-2} 时，这些数才可能非零；在当前情形，
该条件为
\begin{equation}\label{eq:7-5-4}
  \sum_{\nu=2}^n(\nu-1)k_\nu=n+d(n+1)-3.
\end{equation}

各数 $N_d(k_2,\ldots,k_n)$ 是\BookRecordIndex{T0144}格罗莫夫--威滕势的系数。这个生成函数把所有（亏格零）
\BookRecordIndex{T0143}格罗莫夫--威滕不变量的信息汇集在一个表达式中。对 $\mathbb{CP}^n$，它具有如下形式。
把 $H^*(\mathbb{CP}^n)$ 中具有复系数的上同调类写成形式线性组合
\[
  a_t:=t_0+t_1c+t_2c^2+\cdots+t_nc^n,
  \qquad t:=(t_0,\ldots,t_n).
\]
则 $\mathbb{CP}^n$ 的\BookRecordIndex{T0144}格罗莫夫--威滕势是形式幂级数
\begin{equation}\label{eq:7-5-5}
  \Phi^{\mathbb{CP}^n}(s,t_0,\ldots,t_n)
  :=\sum_{d\ge0}\sum_{k\ge0}\frac{s^d}{k!}
    \operatorname{GW}_{dL,k}^{\mathbb{CP}^n}(a_t,\ldots,a_t).
\end{equation}
这里也纳入了 $k<3$ 的项（按注~\ref{rem:7-5-2} 定义），因为这会加强 $\Phi$ 的形式性质。
由命题~\ref{prop:7-5-6} 和注~\ref{rem:7-5-2}，类 $d=0$ 对这个幂级数的贡献是如下三次
多项式：
\begin{align*}
  \phi^{\mathbb{CP}^n}(t)
  &:=\frac16\int_Ma_t\smile a_t\smile a_t\\
  &=\frac16\sum_{i+j+k=n}t_it_jt_k\\
  &=\sum_{i<j<n-i-j}t_it_jt_{n-i-j}
    +\frac12\sum_{\substack{2i\le n\\n\ne3i}}t_i^2t_{n-2i}
    +\frac16t_{n/3}^3.
\end{align*}
这里最后一项只在 $n$ 能被 $3$ 整除时出现。

\begin{lemma}[Kontsevich--Manin]\label{lem:7-5-13}
当 $n\ge2$ 时，
\[
  \Phi^{\mathbb{CP}^n}(s,t)
  =\phi^{\mathbb{CP}^n}(t)
   +\sum_{d>0}\sum_{k_2,\ldots,k_n}
    N_d(k_2,\ldots,k_n)
    \frac{t_2^{k_2}\cdots t_n^{k_n}}{k_2!\cdots k_n!}
    e^{dt_1}s^d.
\]
\end{lemma}

\begin{proof}
省略上标 $\mathbb{CP}^n$。由（零类）公理，
\[
  \phi(t)=\frac16\operatorname{GW}_{0,3}(a_t,a_t,a_t)
  =\sum_{k\ge3}\frac1{k!}\operatorname{GW}_{0,k}(a_t,\ldots,a_t).
\]
此外，由（基本类）公理和习题~\ref{ex:7-5-3}，差
$\Phi_0:=\Phi-\phi$ 与 $t_0$ 无关，并且只含 $k\ge2$ 的 $k$ 点不变量。因此
\begin{align*}
  \Phi_0(s,t)
  &=\sum_{d>0}\sum_{k\ge2}\frac{s^d}{k!}
    \operatorname{GW}_{dL,k}(a_t,\ldots,a_t)\\
  &=\sum_{d>0}\sum_{k\ge2}\sum_{\nu_1,\ldots,\nu_k}
    \frac{s^d}{k!}
    \operatorname{GW}_{dL,k}(c^{\nu_1},\ldots,c^{\nu_k})
    t_{\nu_1}\cdots t_{\nu_k}\\
  &=\sum_{d>0}\sum_{k\ge2}
    \sum_{k_1+\cdots+k_n=k}
    \frac{s^dt_1^{k_1}\cdots t_n^{k_n}}{k_1!\cdots k_n!}
    \operatorname{GW}_{dL,k}
      (c^{\otimes k_1},(c^2)^{\otimes k_2},\ldots,(c^n)^{\otimes k_n})\\
  &=\sum_{d>0}\sum_{k_1+\cdots+k_n\ge2}
    d^{k_1}N_d(k_2,\ldots,k_n)
    \frac{s^dt_1^{k_1}\cdots t_n^{k_n}}{k_1!\cdots k_n!}\\
  &=\sum_{d>0}\sum_{k_2+\cdots+k_n\ge2}
    N_d(k_2,\ldots,k_n)
    \frac{t_2^{k_2}\cdots t_n^{k_n}}{k_2!\cdots k_n!}
    e^{dt_1}s^d.
\end{align*}
第四个等号来自（除子）公理。最后一个等号使用了如下事实：当 $d>0$ 且
$k_2+\cdots+k_n<2$ 时，$N_d(k_2,\ldots,k_n)=0$。这就证明了引理。
\end{proof}

用\BookRecordIndex{T0144}格罗莫夫--威滕势表示，（零类）公理可以写成偏微分方程
\begin{equation}\label{eq:7-5-6}
  \partial_0\partial_i\partial_j\Phi=g_{ij}.
\end{equation}
这里 $\partial_i=\partial/\partial t_i$，常数 $g_{ij}$ 如（分裂）公理中所定义。对
$H^*(\mathbb{CP}^n)$ 的标准基，它们为
\[
  g_{ij}=g^{ij}:=
  \begin{cases}
    1,&i+j=n,\\
    0,&i+j\ne n.
  \end{cases}
\]
此外，事实证明，对任意四个整数
$i,j,k,\ell\in\{0,\ldots,n\}$，函数
$\Phi=\Phi^{\mathbb{CP}^n}$ 都满足如下非凡的二次三阶偏微分方程：
\begin{equation}\label{eq:7-5-7}
  \sum_{\nu,\mu=0}^n
  \partial_i\partial_j\partial_\nu\Phi\,g^{\nu\mu}\,
  \partial_\mu\partial_k\partial_\ell\Phi
  =\sum_{\nu,\mu=0}^n
  \partial_j\partial_k\partial_\nu\Phi\,g^{\nu\mu}\,
  \partial_\mu\partial_i\partial_\ell\Phi.
\end{equation}
这称为 WDVV 方程（见 Dubrovin \srcite{93}）。它是一族高度超定的偏微分方程；令人惊讶的
是，它居然有解。Kontsevich 和 Manin 在 \srcite{216} 中说明，（分裂）公理怎样蕴含
\BookRecordIndex{T0144}格罗莫夫--威滕势满足 WDVV 方程。亏格零\BookRecordIndex{T0143}格罗莫夫--威滕不变量在\BookRecordIndex{T0058}半正情形满足（分裂）公理，
最先由 Ruan--Tian \srcite{345} 证明（稍后 Liu \srcite{248} 以及本书 1994 年版也对单调
情形给出了证明）。这里沿用 Kontsevich--Manin 的论证，说明定理~\ref{thm:7-5-10} 蕴含
$\Phi^{\mathbb{CP}^n}$ 满足式~\eqref{eq:7-5-7}。更详尽的讨论见第~11 章。

\begin{proposition}\label{prop:7-5-14}
$\mathbb{CP}^n$ 的\BookRecordIndex{T0144}格罗莫夫--威滕势 $\Phi^{\mathbb{CP}^n}$ 满足 WDVV 方程。
\end{proposition}

\begin{proof}
我们说明，$\mathbb{CP}^n$ 的 WDVV 方程来自定理~\ref{thm:7-5-10} 中的\BookRecordIndex{T0254}粘合规则。注意，
$\Phi$ 的三阶导数可以写成
\[
  \partial_i\partial_j\partial_\nu\Phi(s,t)
  =\sum_{d\ge0}\sum_{m\ge0}\frac{s^d}{m!}
    \operatorname{GW}_{dL,m+3}
    (a_t,\ldots,a_t,c^i,c^j,c^\nu).
\]
因此，式~\eqref{eq:7-5-7} 左端具有形式
\[
  \sum_{\nu=0}^n
  \partial_i\partial_j\partial_\nu\Phi(s,t)\,
  \partial_{n-\nu}\partial_k\partial_\ell\Phi(s,t)
  =\sum_{d\ge0}\sum_{m\ge0}\frac{s^d}{m!}
    \Psi_{ijk\ell}^{d,m}(t),
\]
其中 $\Psi_{ijk\ell}^{d,m}$ 是 $t$ 的 $m$ 次齐次多项式
\begin{align*}
  \Psi_{ijk\ell}^{d,m}(t)
  &:={}
  \sum_{d_0+d_1=d}\sum_{m_0+m_1=m}\sum_{\nu=0}^n
  \binom{m}{m_0}
  \operatorname{GW}_{d_0L,m_0+3}
    (a_t,\ldots,a_t,c^i,c^j,c^\nu)\\
  &\hspace{8em}\cdot
  \operatorname{GW}_{d_1L,m_1+3}
    (c^{n-\nu},c^k,c^\ell,a_t,\ldots,a_t).
\end{align*}
WDVV 方程断言
\[
  \Psi_{ijk\ell}^{d,m}(t)=\Psi_{jki\ell}^{d,m}(t).
\]
由于全部上同调类的次数都是偶数，这个等式来自定理~\ref{thm:7-5-10} 的第二个恒等式。
\end{proof}

原则上，可以用 WDVV 方程计算 $\mathbb{CP}^n$ 的\BookRecordIndex{T0144}格罗莫夫--威滕势的系数
$N_d(k_2,\ldots,k_n)$。Kontsevich 和 Manin 在 \srcite{216}，定理 3.1 中证明，这些系数
由这个方程以及事实 $N_1(0,\ldots,0,2)=1$ 唯一决定。在定理~\ref{thm:7-4-1} 的证明中，
我们已经
看到
\[
  N_d(2d-2,0,\ldots,0,d+1)\ge N_d.
\]
下面三个习题中，简记
\[
  \Phi_{ijk}:=\partial_i\partial_j\partial_k\Phi,
  \qquad i,j,k\in\{0,1,\ldots,n\}.
\]

\begin{exercise}\label{ex:7-5-15}
设 $n=1$，证明
\[
  \Phi^{\mathbb{CP}^1}=\frac12t_0^2t_1+e^{t_1}s.
\]
这里 $\Phi_{001}=1$，而 WDVV 方程不包含新信息。
\end{exercise}

\begin{exercise}\label{ex:7-5-16}
设 $n=2$，证明
\[
  \Phi^{\mathbb{CP}^2}
  =\frac12t_0^2t_2+\frac12t_0t_1^2
   +\sum_{d>0}N_d\frac{t_2^{3d-1}}{(3d-1)!}e^{dt_1}s^d.
\]
这里 $\Phi_{011}=\Phi_{002}=1$。证明 WDVV 方程等价于
\[
  \Phi_{111}\Phi_{122}+\Phi_{222}=(\Phi_{112})^2,
\]
并验证比较系数会给出命题~\ref{prop:7-5-11} 的递推公式。
\end{exercise}

\begin{exercise}\label{ex:7-5-17}
设 $n=3$。此时式~\eqref{eq:7-5-4} 为 $k_2+2k_3=4d$。证明\BookRecordIndex{T0144}格罗莫夫--威滕势为
\[
  \Phi^{\mathbb{CP}^3}
  =t_0t_1t_2+\frac12t_0^2t_3+\frac16t_1^3
   +\sum_{d=1}^\infty\sum_{k=0}^{2d}
    N_d(4d-2k,k)
    \frac{t_2^{4d-2k}t_3^k}{(4d-2k)!k!}e^{dt_1}s^d.
\]
这里 $\Phi_{012}=\Phi_{003}=1$，而对 $i\le j$，其他所有 $\Phi_{0ij}$ 都为零。由于
$g_{03}=g_{12}=g_{21}=g_{30}=1$，把 $i,j,k,\ell$ 取遍 $\{1,2,3\}$，WDVV 方程给出如下
六个非平凡方程：
\begin{align}
  2\Phi_{123}
    &=\Phi_{111}\Phi_{222}-\Phi_{112}\Phi_{122},\notag\\
  \Phi_{133}
    &=\Phi_{111}\Phi_{223}-\Phi_{113}\Phi_{122},\notag\\
  \Phi_{233}
    &=\Phi_{222}\Phi_{113}-\Phi_{112}\Phi_{223},\notag\\
  \Phi_{333}
    &=(\Phi_{123})^2+\Phi_{113}\Phi_{223}
      -\Phi_{122}\Phi_{133}-\Phi_{112}\Phi_{233},\label{eq:7-5-8}\\
  2\Phi_{123}\Phi_{113}
    &=\Phi_{111}\Phi_{233}+\Phi_{112}\Phi_{133},\notag\\
  2\Phi_{123}\Phi_{223}
    &=\Phi_{222}\Phi_{133}+\Phi_{122}\Phi_{233}.\notag
\end{align}
证明：如果 $i,j,k,\ell$ 中有一个为零，则方程恒真；最后两个方程可由前三个\BookRecordIndex{T0326}推出。因此，
$\mathbb{CP}^3$ 的 WDVV 方程等价于式~\eqref{eq:7-5-8} 中的前四个方程。

比较系数，把 WDVV 方程化成数
\[
  N_{d,k}:=N_d(4d-2k,k),
  \qquad d\ge1,\quad k=0,1,2,\ldots,2d
\]
的递推公式。因此，$N_{d,k}$ 是 $\mathbb{CP}^3$ 中通过 $k$ 个一般点和
$4d-2k$ 条一般直线的 $d$ 次\BookRecordIndex{T0386}有理曲线数。证明式~\eqref{eq:7-5-8} 的第一个方程给出
\begin{align*}
  2dN_{d,k+1}-N_{d,k}
  ={}&\sum_{\substack{d_0+d_1=d\\k_0+k_1=k}}
  \binom{k}{k_0}
  \left(
    d_0^3\binom{4d-2k-3}{4d_0-2k_0}
    -d_0^2d_1\binom{4d-2k-3}{4d_0-2k_0-1}
  \right)\\
  &\hspace{16em}\cdot N_{d_0,k_0}N_{d_1,k_1}
\end{align*}
对 $0\le k\le2d-\frac32$ 成立。证明第二个方程给出
\begin{align*}
  dN_{d,k+1}-N_{d,k}
  ={}&\sum_{\substack{d_0+d_1=d\\k_0+k_1=k}}
  \binom{4d-2k-2}{4d_0-2k_0}
  \left(
    d_0^3\binom{k-1}{k_0}
    -d_0^2d_1\binom{k-1}{k_1}
  \right)\\
  &\hspace{16em}\cdot N_{d_0,k_0}N_{d_1,k_1}
\end{align*}
对 $1\le k\le2d-1$ 成立。证明第三个方程给出
\begin{align*}
  N_{d,k+1}
  ={}&\sum_{\substack{d_0+d_1=d\\k_0+k_1=k}}
  \binom{k-1}{k_0}
  \left(
    d_1^2\binom{4d-2k-3}{4d_1-2k_1}
    -d_0^2\binom{4d-2k-3}{4d_0-2k_0-1}
  \right)\\
  &\hspace{16em}\cdot N_{d_0,k_0}N_{d_1,k_1}
\end{align*}
对 $1\le k\le2d-\frac32$ 成立。这些公式再次表明 WDVV 方程高度超定。把它们与任意两个
不同点之间有唯一一条 $\mathbb{CP}^3$ 中的直线这一事实结合，证明
\[
  N_{1,2}=1,\qquad N_{1,1}=1,\qquad N_{1,0}=2,
\]
\[
  N_{2,4}=0,\qquad N_{2,3}=1,\qquad N_{2,2}=4,
  \qquad N_{2,1}=18,\qquad N_{2,0}=92,
\]
\[
  N_{3,6}=1,\qquad N_{3,5}=5,\qquad N_{3,4}=30,
  \qquad N_{3,3}=190.
\]
考察这些数的几何意义。直接验证 $N_{1,1}=1$ 与 $N_{1,0}=2$（见例~\ref{ex:7-1-14}
和习题~\ref{ex:7-1-16}），并求 $N_{3,2},N_{3,1},N_{3,0}$ 的值。
\end{exercise}
